% =============================================================================
% First-year progression report -- CONDENSED SUMMARY VERSION
% Topological Feature Learning for Parameter Estimation in Neyman-Scott
% Processes -- Flavio Gualtieri
% =============================================================================

\documentclass[11pt,a4paper]{article}

\usepackage[margin=1in]{geometry}
\usepackage{mathpazo}
\usepackage{microtype}

\usepackage{amsmath}
\usepackage{amssymb}
\usepackage{amsthm}

\usepackage{graphicx}
\graphicspath{{../figs/}}
\usepackage{booktabs}
\usepackage{enumitem}
\usepackage[font=small,labelfont=bf,skip=4pt]{caption}
\setlength{\tabcolsep}{5pt}
\renewcommand{\arraystretch}{1.0}
\setlength{\abovedisplayskip}{5pt}
\setlength{\belowdisplayskip}{5pt}
\setlength{\abovedisplayshortskip}{3pt}
\setlength{\belowdisplayshortskip}{3pt}

\usepackage{natbib}
\usepackage[hidelinks]{hyperref}
\raggedbottom
\renewcommand{\textfraction}{0.05}
\renewcommand{\topfraction}{0.95}
\renewcommand{\bottomfraction}{0.95}
\renewcommand{\floatpagefraction}{0.85}

% ----------------------------------------------------------------------------
% Macros
% ----------------------------------------------------------------------------
\newcommand{\R}{\mathbb{R}}
\newcommand{\N}{\mathbb{N}}
\newcommand{\Z}{\mathbb{Z}}
\newcommand{\E}{\mathbb{E}}
\newcommand{\pc}{\mathbf{x}}
\newcommand{\simpcom}{\Sigma}
\newcommand{\filt}{\mathcal{F}}
\newcommand{\kpar}{\kappa}
\newcommand{\mo}{\mu}

% ----------------------------------------------------------------------------
\title{Topological Feature Learning for Parameter Estimation\\
       in Neyman-Scott Processes\\[4pt]
       \large First-year progression report}
\author{Flavio Gualtieri}
\date{\today}

\begin{document}
\maketitle

% #############################################################################
\section{Introduction}
\label{sec:intro}
% #############################################################################

The study of spatial point processes finds applications in astronomy, cell
biology, epidemiology, etc \citep{moller2004statistical}. Many natural phenomena
-- such as galaxy clustering -- are modelled using SPPs, so one can hope to
leverage an understanding of these processes to obtain novel insights. The
parameters of spatial point processes are of particular interest because they
are interpretable scientific quantities. For example, when Neyman-Scott
processes are used to model galaxy clustering the parent intensity $\kappa$
corresponds to the number density of haloes hosting galaxies, the offspring
intensity $\mu$ is the mean occupancy of a halo, and the cluster scale $\sigma$
is the spatial extent of the galaxy distribution within a halo. In forest
ecology, \citet{waagepetersen2009twostep} fit Thomas processes to three tropical
tree species and read the cluster scale directly as a summary of how tightly
seeds are dispersed, finding that species dispersed by exploding capsules, by
wind, and by birds and mammals separate along that parameter.

A number of parameter estimation methods have been proposed, typically relying
on closed-form expressions for summary statistics such as the $K$ and $L$
functions. Minimum contrast estimation fits parameters by numerically minimizing
a discrepancy between the empirical and theoretical curve, most commonly for $K$
or $g$ \citep{diggle2003spatial}. Palm likelihood instead builds a composite
likelihood from the process's Palm intensity, and has been developed
specifically for Neyman-Scott processes \citep{tanaka2008parameter}. Fully
Bayesian approaches treat the unobserved parent locations as latent variables
and sample the posterior over $\theta$ by MCMC \citep{moller2004statistical}.

These methods share three limitations. First, they rely on closed-form summary
statistics, which may not exist (the LGCP-Strauss process of
\citet{vihrs2022neural} has no tractable density) or may exist but fail the
regularity conditions the estimator's theory needs (the two-step estimator of
\citet{waagepetersen2009twostep} requires $K_\psi$ to be twice continuously
differentiable in $\psi$, which excludes the Mat\'ern cluster process). Second,
they are sensitive to hyperparameters chosen per dataset by rule of thumb, with
no held-out quantity indicating whether the choice was good: minimum contrast
depends on $(r_{max}, q, p)$, and \citet{waagepetersen2009twostep} report that
Thomas estimates of $\kappa$ and $\omega$ are ``quite sensitive'' to $r_{max}$
and $q$, being strongly negatively correlated so that the product
$\hat\kappa\hat\omega$ is stable while the individual estimates are not. Third,
they destabilise as the process approaches complete spatial randomness: writing
$\lambda=\kappa\mu$ for the expected number of points per unit area and
$\nu = \sigma\sqrt{\kappa}$ for the ratio of cluster radius to mean parent
spacing, once $\nu \gtrsim 1$ neighbouring clusters overlap and many
$(\kappa,\mu)$ pairs with the same product produce nearly the same observable
pattern. The theoretical $K$-function makes this explicit -- for a Thomas
process $K(t) = \pi t^2 + (1-e^{-t^2/4\sigma^2})/\kappa$, so as $\sigma$ grows
at fixed $t$ the second term vanishes and $K$ tends to the Poisson $K$
regardless of $\kappa$ -- and the contrast surface is close to flat along this
trade-off, so estimates become sensitive to the starting point.

This prompted me to consider whether persistence-based features could be used in
this context, and whether they have advantages over these classical methods:
they provide a multi-scale, process-agnostic descriptor of a point cloud that
requires neither a closed form nor a choice of summary statistic.

% -----------------------------------------------------------------------------
\subsection{Problem setup}
\label{ssec:setup}
% -----------------------------------------------------------------------------

Let $X$ be a Neyman-Scott process on $\R^d$ with unknown parameter vector
$\theta \in \R^{n_{params}}$. For example, the homogeneous Thomas process has
$\theta = (\kappa, \mu, \sigma)$ corresponding to parent intensity, offspring
intensity and cluster scale, respectively. We observe a single realization of
$X$ on a bounded window $W \subset \R^d$, giving a point cloud
$\pc = \{x_1, \dots, x_n\} \subset \R^d$ with $n = |X \cap W|$. The objective is
to produce an estimate $\hat\theta = \hat\theta(\pc)$ of $\theta$ from $\pc$
alone. We stress \emph{single realization}: $\theta$ is not recovered from
repeated draws of $X$. This is the setting a practitioner faces: one observed
pattern and one chance to estimate the generating parameters. It is also what
makes the problem hard in a way no estimator can fix: a fixed $\theta$ already
induces a distribution over point clouds, and it is this irreducible spread --
not any weakness of a particular estimator -- that sets an error floor.

The model is trained on a set of synthetic clouds
$\{\pc_i, \theta_i\}_{i=1}^{N_{train}}$ by backpropagating the RMS loss
    \[
    \mathcal{L}(\hat\theta_i, \theta_i) = \sqrt{\sum_{j=1}^{n_{params}}(\hat\theta^j_i - \theta^j_i)^2}.
    \]
The scale of each parameter varies; therefore -- if not normalized -- the
large-scale parameters will dominate this loss. To address this, we train in
log-normalized parameter space, mapping
$\theta_i \mapsto (\log(\theta_i) - \mu_{\log\theta})/\sigma_{\log\theta}$,
where -- abusing notation -- $\mu_{\log\theta}, \sigma_{\log\theta}$ denote the
mean and standard deviation of the logged parameters over the training set. All
losses reported below are in this space, so that $\mathcal{L}\approx 1$
corresponds to chance level.

% #############################################################################
\section{Background}
\label{sec:background}
% #############################################################################

\citet{vihrs2022neural} first proposed the use of deep learning for parameter
estimation. \citet{yip2025learning} employ a similar architecture but use
persistence images as the input feature, mapping $H_0$--$H_2$ images of an
$\alpha\mathrm{DTM}\ell$ filtration of a halo catalogue to
$(\Omega_m,\sigma_8)$; their targets are cosmological parameters rather than the
parameters of the underlying point process, and they do not systematically
compare topological summaries. Which summary to use is a live question, since
persistence diagrams do not live in a Hilbert space and cannot be fed directly
to a machine learning model: \citet{adams2017persistence} propose persistence
images, \citet{umeda2017time} Betti curves, \citet{bubenik2015landscapes}
persistence landscapes and \citet{chazal2014silhouette} silhouettes, while
\citet{reininghaus2015kernel} avoid vectorization altogether with a kernel on
diagrams.

% -----------------------------------------------------------------------------
\subsection{Neyman-Scott processes}
\label{ssec:spp}
% -----------------------------------------------------------------------------

A spatial point process on $W \subseteq \R^d$ is characterised by its intensity
function $\rho(u)$, which satisfies $\E[N(W)] = \int_W \rho(u)\,du$ for
$N(W) = |W \cap X|$; when $\rho$ is constant the process is homogeneous, the
canonical example being the Poisson process, whose points are scattered
uniformly. A Neyman-Scott process is a clustered process built on top of one, and
the processes used throughout this work are generated by the standard
construction: expand the window $W$ by an edge buffer to give $W_{exp}$, so that
clusters with parents outside $W$ are not truncated; sample a homogeneous
Poisson \emph{parent} process of intensity $\kappa$ on $W_{exp}$; for each
parent $p$ draw an offspring count $n_p \sim \mathrm{Poisson}(\mu)$; place each
offspring at $p + \epsilon$ with displacement $\epsilon$ drawn from a kernel
$k$; finally discard the parents and retain the offspring falling in $W$.

Every Neyman-Scott process therefore has the parent and offspring intensities
$(\kappa,\mu)$ as parameters, and the choice of kernel supplies the rest. The
only thing that changes between the families considered here is that kernel: for
the \emph{Thomas} process $k = \mathcal{N}(0,\sigma^2 I_d)$, giving
$\theta = (\kappa,\mu,\sigma)$; for the Mat\'ern cluster process $k$ is uniform
on a ball of radius $R$; and for the \emph{nested Thomas} process the parent
process at the second step is itself a Thomas process rather than a Poisson one,
producing clusters of clusters and a five-parameter target.

% -----------------------------------------------------------------------------
\subsection{Persistent homology}
\label{ssec:persistence}
% -----------------------------------------------------------------------------

Radially averaged statistics -- such as the $g$ and $K$ functions -- collapse a
point cloud to a real-valued curve indexed by scale $r$, and are turned into
features by evaluating over a range $[r_{min}, r_{max}]$ and subsampling to
obtain a vector in $\R^n$. This discards multi-scale connectivity and void
structure, i.e.\ \emph{how} the density is distributed across space. Persistent
homology (PH) instead tracks a growing family of simplicial complexes built on
the cloud. A \emph{filtration function} $f$ assigns to each subset of points a
scale at which it appears, subject to a monotonicity condition, and the
sublevel sets $\filt_r = \{\sigma : f(\sigma)\le r\}$ form a nested family of
complexes indexed by $r$. The canonical choice is Vietoris-Rips, where a
simplex enters at its diameter; the DTM$_k$ filtration \citep{anai2020dtm}
instead uses each point's root-mean-square distance to its $k$ nearest
neighbours, so a simplex enters according to how densely its neighbourhood is
occupied rather than how close its points happen to be.

As $r$ grows, homological features of the complexes -- connected components
($H_0$), loops ($H_1$), and higher-dimensional voids -- are \emph{born} and
later \emph{die} by merging into older features. Each feature is recorded as a
pair $(b,d)$ of birth and death scales, and the multiset of all such pairs in
dimension $h$ is the \emph{persistence diagram} $D_h$; the lifetime $d-b$
measures how robust a feature is, so points near the diagonal are usually read
as noise. Applied to a point pattern, $H_0$ records the order and scales at
which clusters merge, and $H_1$ records the gaps and voids the clustering leaves
behind, neither of which survives radial averaging. Crucially, a persistence
diagram requires no closed-form expression and no assumption about which family
generated the pattern. Diagrams are multisets rather than vectors, so before
they can be fed to a network they must be \emph{vectorized}; the persistence
image \citep{adams2017persistence} does this by mapping to birth-persistence
coordinates $(b, d-b)$, placing a weighted Gaussian bump at each point, and
integrating the resulting surface over a $64\times64$ pixel grid.

% #############################################################################
\section{ParamNet}
\label{sec:progress}
% #############################################################################

The work completed so far consists of the simulator and its validation, the
ParamNet model and its architectural ablations, comparison with alternative
vectorizations, and benchmarking against existing baselines trained on the same
data.

% -----------------------------------------------------------------------------
\subsection{Synthetic cloud generation}
\label{ssec:simulator}
% -----------------------------------------------------------------------------

Since the estimator is trained rather than derived, the simulator plays the role
that a closed-form summary statistic plays in a classical method: it defines
what the model is asked to learn. Training data are pairs $(\pc_i, \theta_i)$
where $\theta_i$ is drawn from a \emph{design distribution} over parameter space
and $\pc_i$ is a single realization of the corresponding process, generated by
the construction of Section~\ref{ssec:spp}. The design distribution therefore
acts as a prior, and the qualitative regimes it covers determine where the
trained model is usable.

Sampling each parameter uniformly and independently from its own interval turns
out to be inadequate, because the regimes of interest occupy curved subsets of
the parameter grid and most of the mass would land on configurations that are
near-Poisson or degenerate. I instead sample in reparameterized coordinates that
carry the same meaning for every family: the expected intensity
$\lambda = \kappa\mu$, the parent intensity $\kappa$, and the dimensionless
overlap index $\nu = r\sqrt{\kappa}$, where $r$ is the cluster scale of the
family in question ($r=\sigma$ for Thomas, $r=\sqrt{\sigma_1^2+\sigma_2^2}$ for
nested Thomas, with the extra coordinate $\rho = \sigma_2/\sigma_1$). Since
$\nu$ measures cluster radius against mean inter-parent spacing, $\nu \ll 1$
gives well-separated clusters and $\nu \gtrsim 1$ clusters that merge. Draws are
log-uniform, since the parameters span orders of magnitude, and are constrained
so that $\lambda \in [150,800]$ (persistence computation scales badly in the
number of points), $\kappa \in [15,120]$ (enough clusters per window to be
informative, few enough to remain resolvable), and $\nu \in [0.10,0.90]$. The
upper limit on $\nu$ deliberately retains the near-Poisson regime in which
classical estimators are least stable, rather than restricting to
well-separated clusters. Rejection acts on the design and not on realizations,
so it does not distort the conditional law of $\pc$ given $\theta$.

% -----------------------------------------------------------------------------
\subsection{Method: ParamNet}
\label{ssec:method}
% -----------------------------------------------------------------------------

Let $\pc$ denote a point cloud, $v$ its vectorized persistence summary, and
$e \in \R^{E + 1}$ the resulting feature vector, where $E$ denotes the encoder's
embedding dimension and the $+1$ is the appended $\log N(\pc)$. ParamNet
follows the pipeline
    \begin{equation}
    \label{eq:pipeline}
    \pc \xrightarrow{\;\text{Vectorization}\;} v
    \xrightarrow{\;\text{Encoder}\;} e
    \xrightarrow{\;\text{FCNN}\;} \hat\theta.
    \end{equation}
The pipeline shape follows the neural-estimator patterns used in
\citet{vihrs2022neural} and \citet{yip2025learning}; what I have added is the
application to a more challenging process (nested Thomas), an explicit
evaluation of the design choices at each stage, and a comparison of TDA-based
vectorizations, on which there has been no formalized discussion in this
setting.

\paragraph{Filtration.} My starting point was the canonical Rips filtration,
but it is built from inter-point distances alone: it is sensitive to the
sparsest connection between two regions and does not register how densely those
regions are occupied. Since the parameters of a point process act on local
density, DTM$_k$ seemed the more natural object, and it is what the experiments
bore out -- Rips is the worst filtration arm on both processes, by a wide margin
on nested Thomas (Table~\ref{tab:branches}). Within DTM, error grows
monotonically with $k$, consistent with a finer $k$-nearest neighbourhood
retaining more of the local-density signal while a coarser one averages it away.
I also tried fusing several scales to give the encoder fine and coarse views
simultaneously; this helps on nested Thomas but not on Thomas, so the reported
models use $k=5$ for Thomas and fused $k\in\{5,10,15\}$ for nested Thomas. That
this is a per-process choice is a weakness I return to in
Section~\ref{sec:future}.

\paragraph{Encoder, head, and ablations.} For Thomas I compute $H_0$ and $H_1$
persistence images of the DTM$_5$ diagram; for nested Thomas, $H_0$ and $H_1$
images for each of the three DTM$_{\{5,10,15\}}$ filtrations. All images are
$64\times64$. Each is passed through a small convolutional stack: three
$3\times3$ blocks with channel widths $(32,64,128)$, each followed by a ReLU and
$2\times2$ max-pooling with spatial dropout ($p=0.2$) between blocks. Each block
is a CoordConv layer \citep{liu2018coordconv}, concatenating two channels
holding each pixel's normalized coordinates before the convolution: a
persistence image's axes are not arbitrary -- birth encodes local density and
persistence encodes separation -- so a pixel's absolute location, not just its
neighbourhood, is informative. The feature map is reduced by adaptive average
pooling and projected to $E=64$, $\log N(\pc)$ is appended, and the result is
passed to an FCNN head with two hidden layers of width $64$ and $32$ (ReLU,
dropout $p=0.1$). When three scales are used they share one encoder and their
embeddings are concatenated. Every experiment uses AdamW ($\text{lr}=10^{-3}$,
weight decay $3\times10^{-4}$), batch size $32$, up to $500$ epochs with
early-stopping patience $50$.

Each of these choices was ablated one axis at a time against this configuration
(Table~\ref{tab:branches}). A shared encoder beats three independent ones. $H_0$
alone is competitive while $H_1$ alone collapses to chance, so the merge
structure carries most of the signal. Of the vectorizations,
Figure~\ref{fig:vectorization_schematic} shows the four compared on a common
diagram, and none of landscapes, silhouettes or Betti curves improved on
persistence images on either process. Two further arms isolate where the signal
lives: training on $\log N$ alone collapses to the chance level on both
processes, while removing $\log N$ from the full model costs little on nested
Thomas ($0.192\to0.203$) and is merely noisy on Thomas. The persistence images,
not the point count, are doing the work.

\begin{figure}[tbp]
    \centering
    \includegraphics[width=\textwidth]{vectorization_schematic.pdf}
    \caption{One point cloud $\rightarrow$ its $H_0/H_1$ persistence diagrams
    $\rightarrow$ the vectorizations compared in Table~\ref{tab:branches}:
    persistence-image raster, landscape stack, silhouette curve, Betti curve,
    and persistence-statistics summary. Top row: the tight-cluster design
    ($\kappa=50$, $\mu=4$, $\sigma=0.01$); bottom row: the near-Poisson design
    ($\kappa=4$, $\mu=5$, $\sigma=0.5$). Statistics-panel bars are min-max
    normalized per statistic, with raw values printed above each bar, since
    count, total persistence, and entropy live on incomparable scales.}
    \label{fig:vectorization_schematic}
\end{figure}

\begin{table}[tbp]
\centering
\footnotesize
\begin{tabular}{llcc}
\toprule
Axis & Arm & Thomas $\mathcal{L}$ & Nested-Thomas $\mathcal{L}$ \\
\midrule
Filtration & Rips (vec\_multik, landscape) & $0.178 \pm 0.008$ & $0.461 \pm 0.016$ \\
           & DTM$_5$ / DTM$_{15}$ & $0.125 \pm 0.007$ / $0.146 \pm 0.006$ & $0.205 \pm 0.003$ / $0.260 \pm 0.017$ \\
           & DTM$_{\{5,10,15\}}$ (fused) & $0.126 \pm 0.008$ & $0.192 \pm 0.010^{c}$ \\
Homology   & $H_0$ only / $H_1$ only & $0.125 \pm 0.006$ / $0.900 \pm 0.019$ & $0.280 \pm 0.250^\dagger$ / $0.932 \pm 0.028$ \\
Encoder    & shared / independent, concat & $0.125 \pm 0.007$ / $0.134 \pm 0.009$ & $0.201 \pm 0.010$ / $0.222 \pm 0.021$ \\
Vectorization & Persistence image & $0.126 \pm 0.008$ & $0.192 \pm 0.010^{c}$ \\
           & Landscape & $0.132 \pm 0.008$ & $0.250 \pm 0.005$ \\
           & Silhouette ($p=1$) & $0.165 \pm 0.013$ & $0.271 \pm 0.012$ \\
           & Betti curve & $0.196 \pm 0.014$ & $0.290 \pm 0.005$ \\
Features   & $\log N$ removed & $0.217 \pm 0.278^\dagger$ & $0.203 \pm 0.007$ \\
           & Topology removed ($\log N$ only) & $0.893 \pm 0.015$ & $0.948 \pm 0.019$ \\
\bottomrule
\end{tabular}
\caption{Branch summary: every axis measured against the multi-scale,
$H_0{+}H_1$, CoordConv starting configuration, $n=10$ seeds per arm unless
noted. Vectorization rows are the native-encoder path. A dagger ($^\dagger$)
marks an arm on which at least one seed failed to converge and collapsed to the
chance level, so the mean is dominated by the failed seed(s).
$^{c}$$n=3$. The DTM$_5$ / fused rows support the process-specific scale choice
rather than ablating a single model.}
\label{tab:branches}
\end{table}

% -----------------------------------------------------------------------------
\subsection{Estimation performance and baseline comparison}
\label{ssec:results}
% -----------------------------------------------------------------------------

I compare ParamNet against the classical and neural baselines of
Section~\ref{sec:intro}: minimum contrast fit to Ripley's $K$ and, separately,
to the pair correlation function $g$, and \texttt{vihrs} -- my re-implementation
of the neural estimator of \citet{vihrs2022neural}: a 1-D CNN over the
edge-corrected, centred Ripley $L(r)-r$ curve plus $\log N(\pc)$, i.e.\ the
pipeline of Section~\ref{ssec:method} with the persistence-image feature
replaced by a single classical radial summary. All methods are run on the same
clouds and splits, which are made by parameter vector so that no parameter value
appears in both training and evaluation. Chance level is measured by retraining
the default model on randomly permuted labels.

\texttt{vihrs} beats ParamNet on Thomas: its per-parameter losses
($\mathcal{L}^\kappa=0.192$, $\mathcal{L}^\mu=0.122$,
$\mathcal{L}^\sigma=0.029$) are each individually lower than ParamNet's own
($0.201$, $0.134$, $0.040$; Table~\ref{tab:perparam}). So on this single-level
process, the topological pipeline does not beat a hand-built radial summary. On
nested Thomas, ParamNet gains an advantage on the density and count targets
while \texttt{vihrs} maintains an edge on both scale parameters. $L(r)-r$ is
built to resolve a characteristic radius directly, while a density-based
topological feature is not, and it is only on the more complex two-level process
that this trade-off favours the topological pipeline.

Both neural methods dominate minimum contrast on Thomas by roughly an order of
magnitude, but for two different reasons. On $K$ the estimator achieves losses
of $0.365$--$0.399$ for $6/10$ seeds -- close to \texttt{vihrs} -- and
$0.903$--$2.056$ on the remaining $4/10$, which is the near-Poisson flatness
described in Section~\ref{sec:intro}. On $g$ every seed is worse than chance,
indicating systematic bias rather than scattered non-convergence; the likely
cause is my contrast exponent $c=1$, set by analogy with
\citet{diggle2003spatial}'s $K$-based rule of thumb rather than tuned, where
$c=1/4$ is the value validated for aggregated patterns. I report the number but
do not read it as a statement about contrast estimation on $g$ in general.

Finally, parent intensity $\kappa$ is the hardest target and the cluster scale
$\sigma$ the easiest, for every method on both processes and by a wide margin.
This ordering holds regardless of estimator, so it looks like a property of the
estimation problem itself rather than an artifact of any one feature set or
architecture.

\begin{table}[tbp]
\centering
\small
\begin{tabular}{lcc}
\toprule
Method & Thomas $\mathcal{L}$ & Nested-Thomas $\mathcal{L}$ \\
\midrule
ParamNet$^{\ddagger}$    & $0.125 \pm 0.007$ & $0.192 \pm 0.010$ \\
vihrs ($L(r)-r$ CNN)     & $0.115 \pm 0.007$ & $0.217 \pm 0.006$ \\
Minimum contrast, $K$   & $0.855 \pm 0.682^{a}$ & ---$^{d}$ \\
Minimum contrast, $g$   & $6.199 \pm 0.339^{b}$ & ---$^{d}$ \\
\midrule
\emph{Chance-level reference (label shuffle)} & \emph{$1.003 \pm 0.019$} & \emph{$1.006 \pm 0.020$} \\
\bottomrule
\end{tabular}
\caption{Estimator comparison on Thomas and Nested Thomas, $n=10$ seeds, on the
shared test split. $^{\ddagger}$ParamNet is process-specific: DTM$_5$ for
Thomas, fused DTM$_{5,10,15}$ for nested Thomas. $^a$ Minimum contrast on $K$ is
unstable across seeds rather than uniformly bad, and $^b$ minimum contrast on
$g$ is consistently worse than chance; both are discussed in the text.
$^{d}$Not run: no closed-form summary statistic for nested Thomas has yet been
derived, so its absence is an outstanding experiment rather than an
impossibility.}
\label{tab:classical-comparison}
\end{table}

\begin{table}[tbp]
\centering
\small
\textbf{(a) Thomas} ($\kappa$ parent intensity, $\mu$ offspring intensity, $\sigma$ cluster scale)\\[2pt]
\begin{tabular}{lcccc}
\toprule
Parameter & ParamNet & vihrs & Min.\ contrast, $K$ & Min.\ contrast, $g$ \\
\midrule
$\mathcal{L}^\kappa$ & $0.201 \pm 0.010$ & $0.192 \pm 0.013$ & $1.373 \pm 1.087$ & $8.825 \pm 0.485$ \\
$\mathcal{L}^\mu$    & $0.134 \pm 0.007$ & $0.122 \pm 0.007$ & $0.939 \pm 0.758$ & $6.334 \pm 0.370$ \\
$\mathcal{L}^\sigma$ & $0.040 \pm 0.004$ & $0.029 \pm 0.002$ & $0.254 \pm 0.203$ & $3.437 \pm 0.170$ \\
\bottomrule
\end{tabular}

\vspace{10pt}
\textbf{(b) Nested Thomas}\\[2pt]
\begin{tabular}{lcc}
\toprule
Parameter & ParamNet$^{c}$ & vihrs \\
\midrule
parent intensity      & $0.283 \pm 0.023$ & $0.320 \pm 0.010$ \\
meta-offspring        & $0.402 \pm 0.022$ & $0.493 \pm 0.022$ \\
meta-cluster scale    & $0.143 \pm 0.006$ & $0.132 \pm 0.009$ \\
mean offspring        & $0.086 \pm 0.004$ & $0.110 \pm 0.010$ \\
cluster scale         & $0.044 \pm 0.000$ & $0.028 \pm 0.003$ \\
\bottomrule
\end{tabular}
\caption{Per-parameter loss $\mathcal{L}^j$ underlying the aggregate
$\mathcal{L}$ of Table~\ref{tab:classical-comparison} (same methods, same
seeds; ParamNet/nested is $n{=}3$). Minimum contrast has no nested-Thomas
columns for the reason given there.}
\label{tab:perparam}
\end{table}

% #############################################################################
\section{Discussion, limitations and future work}
\label{sec:future}
% #############################################################################

The results support three claims. Persistence images are informative about
Neyman-Scott parameters: both processes sit roughly an order of magnitude below
the label-shuffle ceiling, and the two ``Features'' arms of
Table~\ref{tab:branches} place that signal in the images rather than in
$\log N$. They are not better than a well-chosen radial summary on a
single-level process: \texttt{vihrs} wins on all three Thomas parameters, and on
both scale parameters of nested Thomas. The advantage is confined to the
two-level process and to its density and count targets, which is what one would
expect of a feature built from local density rather than from a characteristic
radius. Two secondary observations: the ordering of parameter difficulty is
stable across feature sets and architectures, so it reads as a property of the
estimation problem; and the margin over minimum contrast is not a single effect,
being seed-to-seed instability on $K$ and a systematic bias traceable to an
untuned contrast exponent on $g$. Neither number is a verdict on contrast
estimation.

\paragraph{Limitations.}
\begin{itemize}[leftmargin=*,itemsep=2pt]
    \item \textbf{Multiple models.} The number of filtration scales is a
    per-process choice with no established rule beyond experimentation and
    comparison of validation losses, which undercuts the process-agnostic
    framing of Section~\ref{sec:intro}. Always fusing would restore a single
    model at $3\times$ the persistence cost and $0.001$ on Thomas, i.e.\ inside
    seed noise; attention over scales might be the fix.
    \item \textbf{Competitiveness.} Three alternative vectorizations were
    implemented and none improved on persistence images. This is evidence about
    these summaries under this architecture and this design distribution, not a
    general claim about vectorizations.
    \item \textbf{No advantage on Thomas.} \texttt{vihrs} wins on every
    parameter of the single-level process, and on both scale parameters of the
    nested one.
    \item \textbf{The design distribution is a prior.} The trained model
    approximates a posterior summary under it, so its ranges must cover the
    intended application domain.
    \item \textbf{Computational cost.} Persistence computation scales badly in
    $n$; this is why the point budget of Section~\ref{ssec:simulator} is capped
    at 800. The inference cost per cloud is currently unmeasured.
\end{itemize}

\paragraph{Planned directions.}
\begin{itemize}[leftmargin=*,itemsep=2pt]
    \item The encoder architecture: aim to design a universal architecture that
    is process-agnostic.
    \item Inhomogeneous / non-stationary processes, where classical closed forms
    are least available and the topological approach has the most to gain. The
    comparator is the two-step procedure of \citet{waagepetersen2009twostep},
    which needs the covariates observed and the intensity correctly specified; a
    simulation-trained estimator needs neither.
    \item Apply to real data. Galaxy catalogues if the running example of
    Section~\ref{sec:intro} is cosmological. The Barro Colorado Island tree data
    of \citet{waagepetersen2009twostep} might be a second target: a small public
    data set already fitted with a Thomas process.
\end{itemize}

% #############################################################################
% REFERENCES
% #############################################################################
{\footnotesize
\setlength{\bibsep}{1pt plus 0.3ex}
\begin{thebibliography}{9}
\setlength{\itemsep}{0pt}
\setlength{\parskip}{0pt}
\bibitem[Adams et al.(2017)]{adams2017persistence}
H.~Adams et al.
\newblock Persistence images: a stable vector representation of persistent
homology.
\newblock \emph{Journal of Machine Learning Research}, 18(8):1--35, 2017.

\bibitem[Anai et al.(2020)]{anai2020dtm}
H.~Anai, F.~Chazal, M.~Glisse, Y.~Ike, H.~Inakoshi, R.~Tinarrage, and
Y.~Umeda.
\newblock DTM-based filtrations.
\newblock In \emph{Topological Data Analysis}, Abel Symposia 15, pages 33--66.
Springer, 2020.

\bibitem[Bubenik(2015)]{bubenik2015landscapes}
P.~Bubenik.
\newblock Statistical topological data analysis using persistence landscapes.
\newblock \emph{Journal of Machine Learning Research}, 16:77--102, 2015.

\bibitem[Chazal et al.(2014)]{chazal2014silhouette}
F.~Chazal, B.~T. Fasy, F.~Lecci, A.~Rinaldo, and L.~Wasserman.
\newblock Stochastic convergence of persistence landscapes and silhouettes.
\newblock In \emph{Proc.\ 30th Annual Symposium on Computational Geometry},
pages 474--483, 2014.

\bibitem[Diggle(2003)]{diggle2003spatial}
P.~J. Diggle.
\newblock \emph{Statistical Analysis of Spatial Point Patterns}.
\newblock Arnold, London, 2nd edition, 2003.

\bibitem[Liu et al.(2018)]{liu2018coordconv}
R.~Liu et al.
\newblock An intriguing failing of convolutional neural networks and the
CoordConv solution.
\newblock In \emph{Advances in Neural Information Processing Systems}, 2018.

\bibitem[M\o ller and Waagepetersen(2004)]{moller2004statistical}
J.~M\o ller and R.~P. Waagepetersen.
\newblock \emph{Statistical Inference and Simulation for Spatial Point
Processes}.
\newblock Chapman \& Hall/CRC, 2004.

\bibitem[Reininghaus et al.(2015)]{reininghaus2015kernel}
J.~Reininghaus, S.~Huber, U.~Bauer, and R.~Kwitt.
\newblock A stable multi-scale kernel for topological machine learning.
\newblock In \emph{IEEE Conference on Computer Vision and Pattern Recognition},
pages 4741--4748, 2015.

\bibitem[Tanaka et al.(2008)]{tanaka2008parameter}
U.~Tanaka, Y.~Ogata, and D.~Stoyan.
\newblock Parameter estimation and model selection for Neyman-Scott point
processes.
\newblock \emph{Biometrical Journal}, 50(1):43--57, 2008.

\bibitem[Umeda(2017)]{umeda2017time}
Y.~Umeda.
\newblock Time series classification via topological data analysis.
\newblock \emph{Transactions of the Japanese Society for Artificial
Intelligence}, 32(3):1--12, 2017.

\bibitem[Vihrs(2022)]{vihrs2022neural}
N.~Vihrs.
\newblock Using neural networks to estimate parameters in spatial point process
models.
\newblock arXiv preprint arXiv:2109.15056, 2022.

\bibitem[Waagepetersen and Guan(2009)]{waagepetersen2009twostep}
R.~Waagepetersen and Y.~Guan.
\newblock Two-step estimation for inhomogeneous spatial point processes.
\newblock \emph{Journal of the Royal Statistical Society: Series B},
71(3):685--702, 2009.

\bibitem[Yip et al.(2025)]{yip2025learning}
J.~H.~T. Yip, A.~Rouhiainen, and G.~Shiu.
\newblock Learning from topology: cosmological parameter estimation from the
large-scale structure.
\newblock \emph{Machine Learning: Science and Technology}, 6:025063, 2025.
\end{thebibliography}
}

\end{document}